\section{Hardware Interface}

The ASUS Xtion Pro RGBD-cameras communicate with the rest of the system via the openni2-package. The nodes of this package opens the camera and transforms the raw data from the camera into various ROS-messages. The RGB-image and the depth-registered point cloud of the main camera are used by the vision node. The depth-registration assigns a 3D point to every pixel in the RGB-image. For pixels where this is not possible, the point is set to (NaN,NaN,NaN). The point cloud from the upward looking camera is sent to the door{\_}detection node. To reduce the computational cost, the frame rates of the cameras are reduced to 5 frames per second, as more images cannot be processed anyway.

The hokuyo{\_}node provides access to the LIDAR. It publishes a laser scan message, which contains an array of range measurements and the information about the angle range and the resolution of the laser scan.

The transformation publisher publishes the transformations of the robot, so the transformation from all the components of the robot into the robot base coordinate frame are available.

The interface to the Turtlebot takes velocity commands, which contain a translational and a rotational velocity. These commands are then executed by the robot, if possible. The interface also provides the data from the odometry as a transformation between the robot and the odometry coordinate frame. 